
% \newpage
\section{Introduction}
Large language models (LLMs) have advanced from basic chatbots \citep{ouyang2022training,achiam2023gpt,hurst2024gpt4o} to sophisticated reasoners \citep{jaech2024openai,guo2025deepseekr1} capable of addressing complex tasks (e.g., Olympiad-level mathematics) through structured problem decomposition, deliberate planning, and reflective self-correction—all without reliance on external knowledge. More recently, the LLM community has begun to explore the development of agentic intelligence \citep{moonshot2025kimik2,zeng2025glm}, which aims to endow LLMs with extended capacities to interface with external tools and environments, 
thereby enabling the autonomous pursuit and completion of more demanding tasks \citep{phan2025humanity}. Along this trajectory, while large-scale pre-training remains the primary driver of model capability \citep{gandhi2025cognitive}, scaling test-time computation is emerging as a complementary avenue for further enhancing LLM intelligence \citep{snell2024scaling}.


As test-time scaling continues to attract growing attention within the LLM community, we posit that two lines of research are poised to become the primary technical drivers of this emerging paradigm. The first is \textit{Reinforcement Learning} (RL) \citep{guo2025deepseekr1,schulman2017proximal}, which offers a principled framework for shaping LLM output distributions via carefully designed reward signals, thereby enabling models to learn effectively from experience \citep{silver2025welcome}. The second, which has received comparatively less attention but may prove even more consequential, is \textit{Task Synthesis} \citep{yang2025swe,li2025websailor,zhao2025promptcot,liu2025webexplorer}—automatic and scalable methods for generating task data that constitute the training ground for RL and thus initiate the process of learning from experience.

\newpage
While human-crafted data continues to play an important role in LLM training \citep{moshkov2025aimo,ahmad2025opencodereasoning}, we argue that synthesized task data will become increasingly central to advancing the frontier of LLM intelligence. The shift is driven by two factors: (1) as LLMs grow more capable, the quality of synthesized data will correspondingly improve; and (2) the tasks posed to LLMs will grow substantially more challenging, with their scope expanding dramatically (e.g., in the pursuit of agentic intelligence). Consequently, reliance on human-generated and curated task data will become prohibitively costly, making scalable synthesis indispensable.

In a broad sense, task synthesis may encompass multiple dimensions, including task definition, problem generation, solution synthesis, environment construction, and even the development of evaluation protocols. In this work, we focus on \textit{Prompt Synthesis} as a simplified yet representative form of task synthesis, where an instance of a task is expressed as a textual prompt\footnote{A prompt may also incorporate multimodal elements such as images, videos, or tools. We defer the exploration of multimodal prompt synthesis to future work.}. Although prompt synthesis has contributed to recent advances in LLMs \citep{yang2025qwen3,zeng2025glm,moonshot2025kimik2,adler2024nemotron}, the underlying technical details are often scarcely disclosed. Publicly available studies, by contrast, typically rely on deliberately engineered prompts to guide powerful LLMs in generating new prompts \citep{wang2023self,xu2024wizardlm,huang2025key,tang2024mathscale,yue2024mammoth2,li2024numinamath,luo2023wizardcoder,wei2023magicoder,huang2024opencoder,xu2025kodcode,liu2025webexplorer,yang2025swe} -an approach that is neither learnable or scalable. As a result, the datasets produced often fall short in difficulty \citep{zhao2025promptcot} and diversity, particularly with respect to expression variation \citep{xu2024wizardlm} and domain coverage (i.e., being restricted to one specific domain). Recently, \citet{zhao2025promptcot} introduces ``rationale''-a form of ``thinking process''-into prompt synthesis, demonstrating substantial improvements in the difficulty of synthesized prompts. This method, however, remains dependent on human-engineered prompts and is confined to the mathematical domain. Consequently, we believe that a fully open-sourced pipeline that integrates advanced prompt synthesis and post-training methods would provide a useful resource for both advancing established capabilities (e.g., mathematical reasoning) and supporting exploration of emerging areas in LLM reasoning (e.g., agentic intelligence, scientific discovery, and beyond).


In this work, we present PromptCoT~2.0, a principled and scalable framework for prompt synthesis that extends the concept–rationale–prompt paradigm beyond the constraints of PromptCoT~1.0~\citep{zhao2025promptcot}.  
Our method introduces an EM-based optimization procedure in which rationales are iteratively refined to guide prompt construction, thereby generating more challenging and diverse problems across both mathematics and programming domains.  
Unlike prior approaches that rely heavily on hand-crafted instructions or domain-specific heuristics, PromptCoT~2.0 is fully learnable and domain-agnostic, enabling synthesis pipelines that can generalize to new reasoning tasks with minimal human intervention.  


We demonstrate the effectiveness of PromptCoT~2.0 through comprehensive experiments under two post-training regimes:  
(1) Self-Play, where synthesized problems provide automatically verifiable feedback that enables models to improve without stronger external teachers; and  
(2) Supervised Fine-Tuning (SFT), where weaker base models are trained from complete reasoning traces distilled from a teacher.  
In the self-play setting, PromptCoT~2.0 delivers substantial improvements over strong Qwen3 baselines and outperforms all existing open-source corpora.  
For example, when applied to \texttt{Qwen3-30B-A3B-Thinking-2507}, our method improves accuracy on AIME~24 from 87.7 to \textbf{92.1}, on AIME~25 from 85.0 to \textbf{89.8}, and on HMMT~Feb~25 from 71.4 to \textbf{76.7}, while also achieving notable gains on LiveCodeBench v5 (+6.1) and v6 (+5.0).  
Overall, these results establish new state-of-the-art performance \emph{at the 30B parameter scale} across all six benchmarks.  
In the SFT setting, we initialize from \texttt{Qwen2.5-7B-Instruct} and train exclusively on prompts synthesized by PromptCoT~2.0.  
Despite relying solely on synthetic problems, the resulting model achieves consistent improvements over both human-curated and hybrid datasets.  
For instance, accuracy improves from 12.8 to \textbf{73.1} on AIME~24, from 8.0 to \textbf{65.6} on AIME~25, from 2.7 to \textbf{46.5} on HMMT~Feb~25, from 14.7 to \textbf{53.4} on LiveCodeBench v5, and from 13.7 to \textbf{48.9} on LiveCodeBench v6, while also achieving the strongest Codeforces performance (706 $\to$ \textbf{1815}). Extended investigation further demonstrates that PromptCoT~2.0 produces prompts that exhibit striking semantic differences from human-created ones and are considerably more challenging than those generated by PromptCoT~1.0.

Our contributions can be summarized as follows:
\begin{itemize}
    \item We introduce a \textbf{new rationale-driven prompt synthesis method}, PromptCoT~2.0, which integrates EM optimization to jointly refine rationale generation and prompt construction. This framework moves beyond handcrafted instructions and scales naturally to multiple domains.
    \item We show how synthesized prompts can be effectively used for \textbf{LLM post-training}. Specifically, we present two complementary regimes—self-play and supervised fine-tuning—and demonstrate consistent improvements over strong Qwen3 and Qwen2.5 baselines across six challenging benchmarks.
    \item We contribute a large-scale corpus of \textbf{synthetic prompts fundamentally different from existing open-source datasets}, with both distributional and difficulty analyses confirming that PromptCoT~2.0 produces problems that are more diverse and substantially more challenging. This resource provides a new foundation for advancing reasoning in open-source LLMs.
\end{itemize}
